\unnumchapter{Обзор литературы}

Каждая подзадача эксплуатации федерации кластеров имеет зрелое решение,
однако ни одно из рассмотренных решений не объединяет политики размещения,
текущие показатели качества работы кластеров и их состояние с точки зрения
аварийного восстановления в одно решение о размещении. Ниже рассматриваются
существующие инструменты по группам задач.

\textbf{Управление федерацией и планирование.} Проект KubeFed~\cite{kubefed}
распространяет ресурсы Kubernetes между кластерами с помощью объектов
\texttt{FederatedResource}, однако не принимает решений на основе метрик, а
его активная разработка фактически прекращена. Admiralty~\cite{admiralty}
размещает поды в удалённых кластерах через виртуальный узел (virtual-kubelet),
но не связан ни с контролем политик, ни с метриками надёжности.
Liqo~\cite{liqo} расширяет кластер ресурсами соседнего через механизм
пиринга и хорошо подходит для сценариев кратковременной нехватки ресурсов, но
не предоставляет детального контроля политик и автоматизации восстановления.
Open Cluster Management~\cite{ocm} предлагает наиболее развитый из
рассмотренных механизм~--- декларативный выбор кластеров по меткам и
ограничениям (Placement API), однако метрики уровня обслуживания в алгоритм
выбора не встроены. Karmada~\cite{karmada} распределяет реплики между
кластерами с учётом их доступной ёмкости. Стандартный планировщик
Kubernetes принимает решения только в пределах одного кластера.

Общим недостатком федеративных планировщиков является выбор кластера по
статическим атрибутам~--- меткам и ресурсной ёмкости~--- без учёта текущей
доли ошибок, задержки и наличия активного инцидента. Как следствие, нагрузка
продолжает направляться в деградирующий кластер до вмешательства человека.

\textbf{Наблюдаемость и SLO.} Стандартом сбора метрик в Kubernetes является
Prometheus~\cite{prometheus}. В федерации у каждого кластера собственный
экземпляр Prometheus, и для расчёта глобальных показателей требуется
агрегация: Thanos~\cite{thanos}, Cortex~\cite{cortex}, Mimir~\cite{mimir} или
VictoriaMetrics~\cite{victoriametrics}. Методика определения целевых
показателей уровня обслуживания и алертинга по скорости расходования бюджета
ошибок описана в~\cite{sre-book, sre-workbook}; инструменты Pyrra~\cite{pyrra}
и Sloth~\cite{sloth} генерируют соответствующие правила из декларативных
описаний, но в границах одного источника данных. Стандарт
OpenTelemetry~\cite{otel} унифицирует трассировку запросов: без сквозного
трейса задержку запроса, проходящего через несколько кластеров, невозможно
разложить на составляющие.

\textbf{Аварийное восстановление и внесение отказов.} Velero~\cite{velero}
выполняет резервное копирование и восстановление состояния кластера, но не
переключает трафик и не принимает решений по метрикам. Argo
Rollouts~\cite{argo-rollouts} и Flagger~\cite{flagger} откатывают выпуск новой
версии при деградации метрик, но работают в границах одного кластера.
Chaos Mesh~\cite{chaos-mesh} и LitmusChaos~\cite{litmus} позволяют вносить
отказы контролируемо и воспроизводимо; без воспроизводимого внесения отказа
сравнение «до» и «после» внедрения механизма не имеет силы.

\textbf{Контроль политик размещения.} Kyverno~\cite{kyverno} (политики на
YAML) и OPA Gatekeeper~\cite{gatekeeper} (язык Rego) работают как admission
webhook внутри кластера и не знают о существовании соседних кластеров.
Политика вида «размещать нагрузку только в кластерах с сертификацией
PCI-DSS» невыразима без данных о других площадках.

\todo{Дополнить обзор научными публикациями: мультикластерное и
геораспределённое планирование, планирование с учётом метрик качества
обслуживания, подходы к измерению надёжности распределённых систем.}

\textbf{Выводы из обзора.} Проведённый анализ выявляет пять узких мест,
каждое из которых в работе закрывается отдельным компонентом контура:
разрозненность инструментов (контур в целом), отсутствие выбора кластера по
текущим метрикам (планировщик уровня федерации), ручная процедура
восстановления (контроллер аварийного восстановления), ограниченность
политик границами одного кластера (двухуровневые политики) и отсутствие
федеративного каталога с инвентарём и состоянием кластеров (каталог
кластеров). Готовые инструменты используются в работе как инфраструктура, а
связывающие их звенья контура разрабатываются.
